

%----------------------------------------------------------------------------------------
%	ENCODING
%----------------------------------------------------------------------------------------

%we use utf8 since we want to build from any machine
\usepackage[utf8]{inputenc}		

%-------------------------------------------------------------------------------
%	LOGIC
%-------------------------------------------------------------------------------


% provides \isempty test
\usepackage{xifthen}

%-------------------------------------------------------------------------------
%	FONT
%-------------------------------------------------------------------------------


% some tex-live fonts - choose your own

% \usepackage[defaultsans]{droidsans}
% \usepackage[default]{comfortaa}
% \usepackage{cmbright}
% \usepackage[default]{raleway}
% \usepackage{fetamont}
\usepackage[default]{gillius}
% \usepackage[light,math]{iwona}
% \usepackage[thin]{roboto} 

% set font default
\renewcommand*\familydefault{\sfdefault} 	
\usepackage[T1]{fontenc}

% more font size definitions
\usepackage{moresize}		

\usepackage{fontawesome}

%-------------------------------------------------------------------------------
%	PAGE LAYOUT  DEFINITIONS
%-------------------------------------------------------------------------------


%debug page outer frames
%\usepackage{showframe}			


%define page styles using geometry
\usepackage[a4paper]{geometry}		

% for example, change the margins to 2 inches all round
\geometry{top=1cm, bottom=-.6cm, left=0.4cm, right=1cm} 	


%less space between header and content
\setlength{\headheight}{-5pt}		


%customize entries left, center and right
%\lhead{}
%\chead{ \small{Jan Küster  $\cdot$ Consultant and Software Engineer $\cdot$  Bremen, Germany  $\cdot$  \textcolor{sectcol}{\textbf{info@jankuester.com}}  $\cdot$ +49 176 313 *** **}}
%\rhead{}


%indentation is zero
\setlength{\parindent}{0mm}

%-------------------------------------------------------------------------------
%	TABLE /ARRAY DEFINITIONS
%-------------------------------------------------------------------------------

%for layouting tables
\usepackage{multicol}			
\usepackage{multirow}

%extended aligning of tabular cells
\usepackage{array}

\newcolumntype{x}[1]{{\raggedleft\hspace{0pt}}p{#1}}%


%-------------------------------------------------------------------------------
%	GRAPHICS DEFINITIONS
%-------------------------------------------------------------------------------

%for header image
\usepackage{graphicx}

%for floating figures
\usepackage{wrapfig}
\usepackage{float}
%\floatstyle{boxed} 
%\restylefloat{figure}

%for drawing graphics		
\usepackage{tikz}				
\usetikzlibrary{shapes, backgrounds,mindmap, trees}


%-------------------------------------------------------------------------------
%	Color DEFINITIONS
%-------------------------------------------------------------------------------
 
\usepackage{transparent}
\usepackage{color}

%accent color
\definecolor{complcol}{RGB}{250,150,10}

%dark background color
\definecolor{bgcol}{RGB}{110,110,110}

%light background / accent color
\definecolor{softcol}{RGB}{225,225,225}

\definecolor{sectcol}{RGB}{0,120,150}


%Package for links, must be the last package used
\usepackage[colorlinks, urlcolor=blue]{hyperref}

%============================================================================%
%
%
%	DEFINITIONS
%
%
%============================================================================%

% returns minipage width minus two times \fboxsep
% to keep padding included in width calculations
\newcommand{\mpwidth}{\linewidth-\fboxsep-\fboxsep}
	

%-------------------------------------------------------------------------------
% 	ARROW GRAPHICS in Tikz
%-------------------------------------------------------------------------------

% a six pointed arrow poiting to the left
\newcommand{\tzlarrow}{(0,0) -- (0.2,0) -- (0.3,0.2) -- (0.2,0.4) -- (0,0.4) -- (0.1,0.2) -- cycle;}	

% include the left arrow into a tikz picture
% param1: fill color
%
\newcommand{\larrow}[1]
{\begin{tikzpicture}[scale=0.58]
	 \filldraw[fill=#1!100,draw=#1!100!black]  \tzlarrow
 \end{tikzpicture}
}

% a six pointed arrow poiting to the right
\newcommand{\tzrarrow}{ (0,0.2) -- (0.1,0) -- (0.3,0) -- (0.2,0.2) -- (0.3,0.4) -- (0.1,0.4) -- cycle;}

% include the right arrow into a tikz picture
% param1: fill color
%
\newcommand{\rarrow}
{\begin{tikzpicture}[scale=0.7]
	\filldraw[fill=sectcol!100,draw=sectcol!100!black] \tzrarrow
 \end{tikzpicture}
}

%-------------------------------------------------------------------------------
%	custom sections
%-------------------------------------------------------------------------------

% create a coloured box with arrow and title as cv section headline
% param 1: section title
%
\newcommand{\cvsection}[1]
{
\colorbox{sectcol}{\mystrut \makebox[1\mpwidth][l]{
\larrow{bgcol} \hspace{-8pt} \larrow{bgcol} \hspace{-8pt} \larrow{bgcol} \textbf{\textcolor{white}{\uppercase{#1}}}\hspace{4pt}
}}\\
}

% create a coloured arrow with title as cv meta section section
% param 1: meta section title
%
\newenvironment{metasection}[1] {
	\vspace{6pt}
	\begin{center}
		\textcolor{white}{\large{\uppercase{#1}}}\\
	\normalsize
	\parbox{0.7\mpwidth}{\textcolor{white}	\hrule}
}{\end{center}}


%-------------------------------------------------------------------------------
%	 CV EVENT
%-------------------------------------------------------------------------------

% creates a stretched box as cv entry headline followed by some paragraphs about 
% the work you did
% param 1:	event time i.e. 2014 or 2011-2014 etc.
% param 2:	event name (what did you do?)
% param 3:	institution (where did you work / study)
% param 4:	list of paragraphs outlining your contributions, see example usage below
%
\newcommand{\cvevent}[4]
{
	\vspace{6pt}
		\begin{tabular*}{1\mpwidth}{p{0.54\mpwidth} >{\raggedleft\arraybackslash}p{0.46\mpwidth}}
			\textcolor{black}{\textbf{#2}} & \textcolor{complcol}{#3}, \textcolor{bgcol}{#1} 
		\end{tabular*}
\vspace{0pt}
\textcolor{softcol}{\hrule}
\vspace{6pt}
	\cvlist{#4}
\vspace{0pt}
}

% Same as \cvevent, but without the institution field and more compact
\newcommand{\cvline}[2]
{
\vspace{4pt}
	\begin{tabular*}{1\mpwidth}{p{0.8\mpwidth} x{0.18\mpwidth}}
 	\textcolor{black}{\cvlist {#2}} & \textcolor{bgcol}{	#1} 
	\end{tabular*}
\vspace{0pt}
\textcolor{softcol}{\hrule}
}

% formats a list of strings with variable length for use in `\cvevent`
% param 1: a list of strings outlining your contributions
\newcommand{\cvlist}[1] {
    \foreach \listitem in {#1}
    {
        \begin{tabular*}
            {1\mpwidth}{p{1\mpwidth}}
            \parbox{1\mpwidth}{\larrow{softcol} \listitem}
            \vspace{2pt}
        \end{tabular*}
    }
}

% creates a stretched box as 
\newcommand{\cveventmeta}[2]
{
	\mbox{\mystrut \hspace{87pt}\textit{#1}}\\
	#2
}

%-------------------------------------------------------------------------------
% CUSTOM STRUT FOR EMPTY BOXES
%-------------------------------------------------------------------------------

\newcommand{\mystrut}{\rule[-.3\baselineskip]{0pt}{\baselineskip}}

%-------------------------------------------------------------------------------
% CUSTOM LOREM IPSUM
%-------------------------------------------------------------------------------
\newcommand{\lorem}
{Lorem ipsum dolor sit amet, consectetur adipiscing elit. Donec a diam lectus.}


% use to vertically center content
% credits to: http://tex.stackexchange.com/questions/7219/how-to-vertically-center-two-images-next-to-each-other
\newcommand{\vcenteredinclude}[1]{\begingroup
\setbox0=\hbox{\includegraphics{#1}}%
\parbox{\wd0}{\box0}\endgroup}

% use to vertically center content
% credits to: http://tex.stackexchange.com/questions/7219/how-to-vertically-center-two-images-next-to-each-other
\newcommand*{\vcenteredhbox}[1]{\begingroup
\setbox0=\hbox{#1}\parbox{\wd0}{\box0}\endgroup}

%-------------------------------------------------------------------------------
%	ICON-SET EMBEDDING
%------------------------------------------------------------------------------ 

% at this point we simplify our icon-embedding by simply referring to a set of png images.
% if you find a good way of including svg without conflicting with other packages you can
% replace this part
\newcommand{\icon}[3]{
	\makebox(#2, #2){\textcolor{#3}{\csname fa#1\endcsname}}
}	%icon shortcut
\newcommand{\icontext}[4]{ 						%icon with text shortcut
	\vcenteredhbox{\icon{#1}{#2}{#4}} \vcenteredhbox{\textcolor{#4}{#3}}
}
\newcommand{\iconhref}[5]{ 						%icon with website url
    \vcenteredhbox{\icon{#1}{#2}{#5}} \href{#4}{\textcolor{#5}{#3}}
}

\newcommand{\iconemail}[5]{ 						%icon with email link
    \vcenteredhbox{\icon{#1}{#2}{#5}} \href{mailto:#4}{\textcolor{#5}{#3}}
}

\newcommand{\customicon}[3]{
	\makebox(#2, #2){\textcolor{#3}{\csname \includegraphics*[width=#2pt, height=#2pt]{#1}\endcsname}}
}	% Custom icon shortcut (#1 is the icon path)

\newcommand{\customicontext}[4]{ %icon with text shortcut
	\vcenteredhbox{\customicon{#1}{#2}{#4}} \vcenteredhbox{\textcolor{#4}{#3}}
}